\section*{Bab \ref{ch:g12:plant-rooted-life} --- Hidup Tumbuhan: Berjaya sambil Berakar}

\begin{solution}{exo:g12:plant-rooted-life:1}
Sistem akar, yang \omterm{prop:g12:plant-rooted-life:surfaces}{rambut akarnya} menyodorkan permukaan luas kepada
tanah, dan sistem tajuk, yang daunnya menyodorkan permukaan luas kepada
cahaya dan udara.
\end{solution}

\begin{solution}{exo:g12:plant-rooted-life:2}
Sebuah liang pada kulit daun yang dibatasi dua \omterm{def:g10:cells-common-unit:cell}{sel} penjaga. Karbon
dioksida masuk, uap air (dan oksigen) keluar. Terbuka dalam cahaya,
tertutup pada malam hari dan dalam kekeringan.
\end{solution}

\begin{solution}{exo:g12:plant-rooted-life:3}
Getah mentah: air dan ion mineral, di dalam \omterm{prop:g12:plant-rooted-life:saps}{xilem}, dari akar ke daun,
ditarik oleh transpirasi. Getah olahan: air dan gula, di dalam \omterm{prop:g12:plant-rooted-life:saps}{floem},
dari daun ke organ yang memakai atau menyimpan gula, ke kedua arah,
didorong oleh tekanan \omterm{def:g10:cells-common-unit:cell}{sel} yang bermuatan gula.
\end{solution}

\begin{solution}{exo:g12:plant-rooted-life:4}
Struktural: kutikula dan kulit kayu, duri (juga rambut penyengat,
silika, \omterm{def:g10:cells-common-unit:organelle}{dinding sel}). Kimiawi: tanin, alkaloid (juga racun, damar,
lateks).
\end{solution}

\begin{solution}{exo:g12:plant-rooted-life:5}
Tanggapan pertumbuhan yang berarah terhadap rangsangan (cahaya,
gravitasi, air). Auksin, yang dibuat di ujung tunas yang tumbuh.
\end{solution}

\begin{solution}{exo:g12:plant-rooted-life:6}
$\qty{40}{cm^2} = \qty{4000}{mm^2}$: $4000 \times 200 = \num{8e5}$
\omterm{def:g12:plant-rooted-life:stomata}{stoma} per helai daun; $\num{1.6e11}$ pada seluruh pohon.
\end{solution}

\begin{solution}{exo:g12:plant-rooted-life:7}
Sekitar tengah hari sampai awal sore. Pada panasnya sore, \omterm{def:g10:cells-common-unit:cell}{sel} penjaga
menutup \omterm{def:g12:plant-rooted-life:stomata}{stomanya} sebagian untuk membatasi kehilangan air, dan
transpirasi mengikuti bukaannya, bukan cahayanya.
\end{solution}

\begin{solution}{exo:g12:plant-rooted-life:8}
Daunnya tetap hijau dan terpasok air (\omterm{prop:g12:plant-rooted-life:saps}{xilemnya}, di dalam kayu, masih
utuh). Akarnya, yang terputus dari gula daun, menghabiskan cadangannya
selama beberapa bulan lalu mati, dan pohonnya ikut mati. Tepat di atas
cincin itu kulit kayunya membengkak oleh gula yang tak dapat lagi
lewat.
\end{solution}

\begin{solution}{exo:g12:plant-rooted-life:9}
Utuh: tanggapannya memang ada. Ujung ditudungi: ujungnyalah tempat
cahaya dirasakan. Pangkal ditudungi: pangkalnya tak perlu disinari
untuk membengkok --- ia menanggapi sebuah isyarat, bukan cahaya. Ujung
dipotong: ujungnya memang diperlukan; isyaratnya datang dari sana.
Bersama-sama keempatnya memisahkan penginderaan (ujung) dari tanggapan
(pangkal) dan menunjukkan bahwa ada sesuatu yang berjalan di antara
keduanya.
\end{solution}

\begin{solution}{exo:g12:plant-rooted-life:10}
Dari matahari: penguapan air di permukaan daun, yang digerakkan panas
surya dan udara kering, menarik kolom air yang malar naik di dalam
\omterm{prop:g12:plant-rooted-life:saps}{xilem}. Tumbuhan itu memasok tabungnya dan \omterm{def:g12:plant-rooted-life:stomata}{stoma} yang terbuka, bukan
energinya.
\end{solution}

\begin{solution}{exo:g12:plant-rooted-life:11}
Bunga angin: kecil, pudar, tanpa harum, tanpa nektar, dengan serbuk
sari ringan berlimpah dan kepala putik berbulu --- angin itu buta dan
gratis, sehingga tumbuhannya menanam modal pada jumlah. Bunga
serangga: besar, berwarna, harum, bernektar, dengan serbuk sari
lengket dalam jumlah sedang --- pengantarnya harus dipikat dan dibayar,
dan ia mengantarkan serbuk sari dengan tepat.
\end{solution}

\begin{solution}{exo:g12:plant-rooted-life:12}
Ia memperoleh air, sehingga \omterm{def:g10:cells-common-unit:cell}{selnya} tetap tegang dan hidup; ia
kehilangan masuknya karbon dioksida, jadi juga \omterm{def:g10:cell-metabolism:autotrophy}{fotosintesis}, serta
pendinginan oleh transpirasi. \omterm{def:g12:plant-rooted-life:stomata}{Stoma} yang tertutup berarti tak ada gula
yang dibuat: tumbuhan itu hidup dari cadangannya, dan kekeringan yang
panjang membuatnya kelaparan sebelum membuatnya kering.
\end{solution}

\begin{solution}{exo:g12:plant-rooted-life:13}
Di dalam daun, ia sebuah pertahanan: racun terhadap pemakan tumbuhan,
yang terpasang di dalam jaringannya. Di dalam nektar, ia sebuah
perekrutan: dosis rendah yang membuat penyerbuknya ingat dan kembali,
dibayarkan bersama gulanya. Satu zat, dua kebutuhan, yang didosis
menurut tempatnya.
\end{solution}

\begin{solution}{exo:g12:plant-rooted-life:14}
Bunga dengan tabung yang sedikit lebih panjang hanya diserbuki ngengat
yang lidahnya mencapai dasarnya, sehingga serbuk sarinya sampai ke
bunga sejenisnya; ngengat dengan lidah yang sedikit lebih panjang
memperoleh nektar yang tak terjangkau yang lain; setiap keuntungan
kecil terseleksi, dan tabung serta lidah memanjang bersama sepanjang
banyak generasi. Bila ngengatnya lenyap, bunganya tak berbiji; bila
bunganya lenyap, ngengatnya kehilangan satu-satunya makanannya.
\end{solution}

\begin{solution}{exo:g12:plant-rooted-life:15}
Tumbuhan mengindera cahaya dengan ujung tunasnya dan gravitasi dengan
akarnya; ia menanggapi dengan tumbuh ke arah yang satu dan menyusuri
yang lain; ia membela diri dengan kulit kayu dan duri, serta membuat
racun dan penghambat pencernaan dalam beberapa jam setelah serangan;
dan ia merekrut serangga untuk membawa serbuk sarinya, burung untuk
membawa bijinya dan tabuhan untuk membunuh ulatnya, dengan bayaran
nektar, buah dan harum. Ia bertindak tanpa berpindah.
\end{solution}

\begin{solution}{pb:g12:plant-rooted-life:1}
\textbf{1.} $\num{250000} \times \qty{25}{cm^2} = \num{6.25e6}\,
\unit{cm^2} = \qty{625}{m^2}$.

\textbf{2.} $\qty{625}{m^2} = \num{6.25e8}\,\unit{mm^2}$; $\times 250
\approx \num{1.6e11}$ \omterm{def:g12:plant-rooted-life:stomata}{stoma}.

\textbf{3.} $625/150 \approx 4$ meter persegi daun per meter persegi
tanah.

\textbf{4.} Sekitar \qty{1000}{m^2} \omterm{prop:g12:plant-rooted-life:surfaces}{rambut akar} berbanding
\qty{625}{m^2} daun: sebanding. Keduanya adalah permukaan pertukaran
yang berukuran menurut aliran yang sama --- air yang hilang dari daun
harus diserap oleh akarnya.

\textbf{5.} Cahaya terserap dalam jaringan setebal sepersekian
milimeter, dan karbon dioksida hanya berdifusi sejauh jarak pendek:
organ yang tebal akan meninggalkan sebagian besar \omterm{def:g10:cells-common-unit:cell}{selnya} dalam gelap
dan kelaparan. Daun tipis yang bertumpuk menempatkan setiap \omterm{def:g10:cells-common-unit:cell}{sel} dekat
cahaya dan udara.

\textbf{6.} $625 \times 0.6 = \qty{375}{L}$.

\textbf{7.} $625 \times 8 = \qty{5}{kg}$ karbon dioksida, yang memberi
sekitar $5/1.47 \approx \qty{3.4}{kg}$ gula (atau ringkasnya: gula yang
dibuat dari \qty{5}{kg} $\mathrm{CO_2}$ sekitar \qty{3.4}{kg}); air
yang dipakai $3.4/1.6 \approx \qty{2.1}{kg}$: sekitar 0.6\% dari
\qty{375}{L} yang ditranspirasikan.

\textbf{8.} Karbon dioksida masuk lewat \omterm{def:g12:plant-rooted-life:stomata}{stoma} terbuka yang sama yang
melepaskan airnya; menutupnya menghentikan \omterm{def:g10:cell-metabolism:autotrophy}{fotosintesis}. Air itulah
harga gulanya.

\textbf{9.} Matahari, yang menguapkan air di daunnya. Mengangkat
\qty{375}{kg} setinggi \qty{20}{m} kira-kira \qty{75}{kJ} kerja ---
sedikit saja, tetapi pompa biologis akan menuntut jaringan, ATP dan
perawatan yang tak perlu dibangun pohon itu.

\textbf{10.} Ia telah menutup sebagian \omterm{def:g12:plant-rooted-life:stomata}{stomanya} terhadap panas yang
kering; masuknya karbon dioksida dan produksi gula turun sebanding ---
sekitar \qty{2}{kg} gula hilang sepanjang sore itu.

\textbf{11.} Sekitar \qty{3.4}{kg} dibuat; \qty{2}{kg} (60\%) ke akar,
batang dan tunas, \qty{0.5}{kg} (15\%) ke bijinya.

\textbf{12.} \omterm{prop:g12:plant-rooted-life:saps}{Floem}, dari daun (sumber) ke akar, batang, tunas dan biji
(lubuk).

\textbf{13.} Disimpan sebagai pati di dalam batang dan akar; musim semi
berikutnya pati itu diubah kembali menjadi gula untuk membangun daun
baru sebelum daunnya dapat memberi makan dirinya sendiri.

\textbf{14.} Gula masih mencapai tunas dan biji di atas cincin itu;
akar dan batang bawah tidak menerima apa-apa: jatah \qty{2}{kg} mereka
turun menjadi sebanyak yang tersimpan saja, dan keduanya kelaparan
sepanjang musim.

\textbf{15.} Daun yang tersisa terus membuat gula, pati yang tersimpan
di batang dan akar menutup kekurangannya, dan pohon itu dapat
mengeluarkan daun baru dari banyak titik tumbuhnya; bijinya adalah
lubuk yang masih dapat diisi pohon itu bila gula musim panasnya cukup.

\textbf{16.} Bunga kecil kehijauan tanpa mahkota, harum atau nektar,
yang menggantung dalam untaian dan menaburkan awan serbuk sari ringan;
kepala putik besar berbulu untuk menangkapnya.

\textbf{17.} Biji yang jatuh di bawah induknya tumbuh, kalaupun
tumbuh, dalam naungannya dan bersaing dengannya; seekor burung jay
membawanya ratusan meter dan menanamnya pada kedalaman yang tepat.
Kehilangan sebagian besar bijinya adalah harga untuk menempatkan
beberapa dengan baik.

\textbf{18.} $3000 \times 0.05 = 150$ biji yang terlupakan; $150/50 = 3$
anakan pohon.

\textbf{19.} Tiga anakan setahun sepanjang beberapa dasawarsa panen
yang baik memberi mungkin seratus anakan sepanjang hidup pohon itu,
dan satu di antaranya harus bertahan sampai dewasa. Angkanya besar
karena hampir setiap biji dan anakan hilang.

\textbf{20.} \qty{625}{m^2} daun dan sekitar $\num{1.6e11}$ \omterm{def:g12:plant-rooted-life:stomata}{stoma};
sekitar \qty{375}{L} diangkat dalam sehari untuk \qty{3.4}{kg} gula;
angin bagi serbuk sarinya dan burung jay bagi bijinya.
\end{solution}
